\section{Μελλοντική Έρευνα και Προοπτικές}
\label{sec:conclusion_future}

Οι προτεινόμενοι άξονες για μελλοντική διερεύνηση περιλαμβάνουν:
\begin{enumerate}
    \item \textbf{Πολυ-αντικειμενική Βελτιστοποίηση (Multi-Objective Pareto):} Επέκταση του αλγορίθμου με NSGA-II/NSGA-III για ταυτόχρονη μεγιστοποίηση του λόγου $L/D$ και ελαχιστοποίηση του δομικού βάρους της ατράκτου.

    Τυπικά, η πολυ-αντικειμενική βελτιστοποίηση ορίζεται μέσω της σχέσης κυριαρχίας Pareto: ένα διάνυσμα $\bm{x}_a$ \emph{κυριαρχεί} (dominates) το $\bm{x}_b$ ως προς ένα σύνολο στόχων $\bm{f} = (f_1, \ldots, f_k)$ αν
    \begin{equation}
    \label{eq:pareto_dominance}
        \forall i \in \{1,\ldots,k\}: f_i(\bm{x}_a) \ge f_i(\bm{x}_b)
        \quad \text{και} \quad
        \exists j: f_j(\bm{x}_a) > f_j(\bm{x}_b).
    \end{equation}
    Το μέτωπο Pareto $\mathcal{F}^*$ αποτελείται από τα μη-κυριαρχούμενα σχέδια. Η ενσωμάτωση NSGA-III στο island model θα επέτρεπε διατήρηση τριτοβάθμιας ποικιλομορφίας (reference-point based diversity) σε κάθε νήσο, ενώ η μετανάστευση ελίτ θα μεταφέρει άτομα κοντά στο μέτωπο Pareto μεταξύ νήσων, δυνητικά επιταχύνοντας τη σύγκλιση στο βέλτιστο μέτωπο.

    \item \textbf{Υβριδική Προσομοίωση VLM + RANS:} Σταδιακή ενσωμάτωση επιλεκτικών προσομοιώσεων Navier-Stokes υψηλής πιστότητας αποκλειστικά για τα άτομα του μετώπου Pareto.

    Ενώ η VLM παρέχει επαρκή ακρίβεια για τα πρώτα στάδια βελτιστοποίησης, η τελική επικύρωση των βέλτιστων σχεδίων απαιτεί επίλυση Navier-Stokes (RANS) για την ακριβή αποτίμηση οριακού στρώματος και φαινομένων αποκόλλησης. Η σταδιακή ενεργοποίηση RANS αξιολογήσεων μόνο για τα υποψήφια άτομα του μετώπου Pareto σε τελικές γενεές θα μείωνε τον συνολικό αριθμό ακριβών αξιολογήσεων, ενώ θα διατηρούσε υψηλή πιστότητα στις τελικές λύσεις.

    \item \textbf{Πειραματική Επικύρωση σε Αεροδυναμική Σήραγγα:} Κατασκευή και δοκιμή των βέλτιστων 3D-εκτυπωμένων πτερύγων σε υποηχητική αεροδυναμική σήραγγα σε συνεργασία με το Τμήμα Μηχανολόγων και Αεροναυπηγών Μηχανικών (MEAD) του Πανεπιστημίου Πατρών. Η σύγκριση μετρηθέντων και προβλεφθέντων τιμών $L/D$ θα επαληθεύσει την εγκυρότητα της ψηφιακής αλυσίδας βελτιστοποίησης.

    \item \textbf{Προσαρμοστικό Surrogate με Bayesian Βελτιστοποίηση:} Αντικατάσταση της στατικής πύλης αβεβαιότητας $\sigma(\bm{x}^*) < \theta$ με ενεργό στρατηγική απόκτησης δεδομένων βασισμένη στο κριτήριο Expected Improvement (EI):
    \begin{equation}
    \label{eq:ei_acquisition}
        \text{EI}(\bm{x}) = \mathbb{E}\bigl[\max(f(\bm{x}) - f^+, 0)\bigr],
    \end{equation}
    όπου $f^+$ είναι η μέχρι τώρα βέλτιστη γνωστή τιμή \cite{jones1998ego}. Ο εκάστοτε κόμβος Surrogate θα επέλεγε τα επόμενα σχέδια για VLM αξιολόγηση βάσει μεγιστοποίησης του EI, μετατρέποντας το σύστημα σε ένα πλήρες Bayesian Optimization loop εντός του εξελικτικού πλαισίου. Αυτή η προσέγγιση έχει αποδειχθεί ιδιαίτερα αποδοτική σε προβλήματα αεροδυναμικής βελτιστοποίησης \cite{forrester2008engineering}.

    \item \textbf{Ομόσπονδο DHT σε Φυσικά Κατανεμημένο Νέφος:} Επέκταση της LEAD DHT τοπολογίας σε ετερογενή υποδομή που εκτείνεται σε πολλαπλούς φυσικούς κόμβους ή νεφοϋπολογιστικές περιοχές (cloud regions). Ο κατανεμημένος δακτύλιος θα ενσωμάτωνε μηχανισμό ασυνεπούς αντιγραφής (eventual consistency) για ανοχή σε κόμβους που αποσυνδέονται, εξασφαλίζοντας συνεχή διαθεσιμότητα της χωρικής ευρετηρίασης. Αυτό θα επέτρεπε κλιμάκωση του συστήματος σε δεκάδες παράλληλες νήσους που εκτελούνται σε γεωγραφικώς κατανεμημένη υποδομή.

    \item \textbf{Ενισχυτική Μάθηση για Προσαρμοστική Πολιτική Μετανάστευσης:} Αντί για σταθερές παραμέτρους $T_{\text{mig}}$ και $K_{\text{mig}}$, εκπαίδευση ενός πράκτορα ενισχυτικής μάθησης (Reinforcement Learning agent) που ρυθμίζει δυναμικά τη συχνότητα και το μέγεθος μετανάστευσης ανάλογα με την παρατηρούμενη ποικιλομορφία κάθε νήσου και τον ρυθμό βελτίωσης καταλληλότητας. Η κατάσταση του πράκτορα θα κωδικοποιούσε στατιστικά στοιχεία πληθυσμού (entropy γονιδίων, μέση τιμή, τυπική απόκλιση καταλληλότητας), και η ανταμοιβή θα ορίζεται από το ποσοστό βελτίωσης καταλληλότητας ανά γενεά, οδηγώντας σε αυτο-ρυθμιζόμενη μετανάστευση που προσαρμόζεται στη δυναμική του εξελικτικού τοπίου.
\end{enumerate}

Συνολικά, το αρχιτεκτονικό πλαίσιο που αναπτύχθηκε στην παρούσα εργασία αποτελεί μια επεκτάσιμη βάση για την υλοποίηση όλων των παραπάνω κατευθύνσεων, δεδομένου ότι η αρθρωτή σχεδίαση μέσω gRPC επιτρέπει την ανεξάρτητη αναβάθμιση κάθε υποσυστήματος χωρίς διακοπή της υπόλοιπης λειτουργίας.
